\documentclass[11pt]{article}
\usepackage{amsmath,amssymb,amsthm,mathtools}
\usepackage[margin=1in]{geometry}

\title{BEpTy: A Calculus of Structured Absence}
\author{Inacio Vasquez}
\date{\today}

\newtheorem{definition}{Definition}
\newtheorem{theorem}{Theorem}
\newtheorem{corollary}{Corollary}

\begin{document}
\maketitle

\section{Primitive datum}

Let \(A \subseteq B \subseteq X\).

\begin{definition}[BEpTy absence object]
\(\mathbb E(A,B) := B \setminus A.\)
\end{definition}

\begin{definition}[Valuation]
For an admissible valuation \(\Phi\), define
\(\beta_\Phi(A,B) := \Phi(B \setminus A).\)
\end{definition}

\section{Core laws}

\begin{theorem}[Zero Law]
\(\mathbb E(A,A)=\varnothing.\)
\end{theorem}

\begin{theorem}[Decomposition Law]
If \(A \subseteq B \subseteq C\), then
\(\mathbb E(A,C)=\mathbb E(A,B)\sqcup \mathbb E(B,C).\)
\end{theorem}

\begin{corollary}[Additive valuation law]
If \(\Phi(\varnothing)=0\) and \(\Phi\) is additive on disjoint unions, then
\(\beta_\Phi(A,C)=\beta_\Phi(A,B)\oplus \beta_\Phi(B,C).\)
\end{corollary}

\section{Finite specialization}

Assume \(X\) is finite and \(\Phi(E)=|E|\). Then
\[
\beta(A,B):=|B\setminus A|.
\]

\begin{theorem}[Finite Detection]
\(\beta(A,B)=0 \iff A=B.\)
\end{theorem}

\begin{theorem}[Finite Unit Classification]
\(\beta(A,B)=1 \iff \exists !x\in X\ \text{such that}\ B=A\cup\{x\}.\)
\end{theorem}

\section{Topological specialization}

Let
\[
\beta_{\mathrm{int}}(A,B):=
\begin{cases}
0,&\operatorname{int}(B\setminus A)=\varnothing,\\
1,&\operatorname{int}(B\setminus A)\neq\varnothing.
\end{cases}
\]

\begin{theorem}[Open Detection]
If \(A \subseteq B \subseteq X\) and \(B\setminus A\) is open in \(X\), then
\(\beta_{\mathrm{int}}(A,B)=0 \iff A=B.\)
\end{theorem}

\begin{corollary}[Closed Shell Detection]
If \(A\) is closed, \(B\) is open, and \(A\subseteq B\), then
\(\beta_{\mathrm{int}}(A,B)=0 \iff A=B.\)
\end{corollary}

\section{Worked interval example}

Let \(A=[0,1]\) and \(B=(0,2)\). Then
\[
B\setminus A=(1,2),
\qquad
\operatorname{int}(B\setminus A)=(1,2)\neq\varnothing,
\qquad
\beta_{\mathrm{int}}(A,B)=1.
\]

\section{Program threshold}

The next threshold is:
\begin{enumerate}
\item prove one nontrivial theorem beyond finite normalization;
\item identify one URF invariant expressible as a BEpTy valuation.
\end{enumerate}

\end{document}
